\setcounter{chapter}{3}
\chapter{Lebesgue内測度}\label{chap:lebesgue-inner}

\paragraph{本章の目標}
Lebesgue内測度，4つの測度の大小関係

前章ではLebesgue外測度 $m^*$ を定義し，
あらゆる集合に対して「外からの体積評価」を与えた．
しかし $m^*$ 単体では加法性が成り立たず，
測度としては不十分であった．

Riemann積分やJordan測度では，
上からの評価（上和，外測度）と
下からの評価（下和，内測度）が一致するとき
「積分可能」あるいは「可測」と定義した．
Lebesgue理論でも同じ戦略を取る：
外測度 $m^*(A)$ と内測度 $m_*(A)$ を比較し，
両者が一致する集合を「Lebesgue可測」と定義する．

本章ではLebesgue内測度を定義し，
外測度との関係およびJordan内測度との関係を調べる．

% ==================================================================
\section{Lebesgue内測度の定義}\label{sec:lebesgue-inner-def}
% ==================================================================

\begin{definition}[Lebesgue内測度]\label{def:lebesgue-inner-measure}
有界集合 $A \subset \RR^d$ に対し，
$A$ を含む有界閉区間 $I \supset A$ を取り，
\[
m_*(A) := |I| - m^*(I \setminus A)
\]
を $A$ の\textbf{Lebesgue内測度}と定める．
\end{definition}

この定義が $I$ の取り方に依らないことを確かめる．

\begin{proposition}[well-definedness]\label{prop:inner-well-defined}
上の定義は $I$ の選び方に依存しない．
\end{proposition}
\begin{proof}
$A \subset I \subset J$ を満たす
有界閉区間 $I, J$ を取る．
$J \setminus A = (J \setminus I) \sqcup (I \setminus A)$
であり，$J \setminus I$ は基本集合だから
補題\ref{lem:outer-additive-elementary}（第3章）より
\[
m^*(J \setminus A)
= m^*(J \setminus I) + m^*(I \setminus A)
= (|J| - |I|) + m^*(I \setminus A).
\]
したがって
\[
|J| - m^*(J \setminus A)
= |J| - (|J| - |I|) - m^*(I \setminus A)
= |I| - m^*(I \setminus A).
\]
すなわち $I$ を $J$ に取り替えても値は変わらない．
任意の二つの区間 $I_1, I_2 \supset A$ に対しては，
共通の区間 $J \supset I_1 \cup I_2$ を取って
$I_1, I_2$ の両方について上の議論を適用すればよい．
\end{proof}

\begin{remark}
Jordan内測度
$m_{J,*}(A) = \sup\{m_J(E) : E \subset A,\ E \text{ は基本集合}\}$
が「$A$ の中に入る基本集合の体積の上限」であったのに対し，
Lebesgue内測度は「$I$ の体積から $A$ の補集合の外測度を引いたもの」である．
表現は異なるが，
いずれも「$A$ の内側からの体積評価」を与えるという動機は共通している．
\end{remark}

\begin{remark}[非有界集合への拡張]\label{rem:inner-unbounded}
非有界集合 $A \subset \RR^d$ に対しては
\[
m_*(A) := \sup\{m_*(A \cap I) : I \text{ は有界閉区間}\}
\]
と定義すれば，Lebesgue内測度を $\calP(\RR^d)$ 全体に拡張できる．
以下では主に有界集合を扱う．
\end{remark}

% ==================================================================
\section{基本性質}\label{sec:lebesgue-inner-prop}
% ==================================================================

\begin{theorem}[Lebesgue内測度の性質]\label{thm:lebesgue-inner-properties}
有界集合 $A, B \subset \RR^d$ に対して次が成り立つ．
\begin{enumerate}
\item （非負性）$m_*(A) \ge 0$．$m_*(\emptyset) = 0$．
\item （内・外の関係）$m_*(A) \le m^*(A)$．
\item （単調性）$A \subset B \Rightarrow m_*(A) \le m_*(B)$．
\item （補集合との関係）$A \subset I$（有界閉区間）ならば
$m_*(A) + m^*(I \setminus A) = |I|$．
\item （平行移動不変性）$m_*(A + c) = m_*(A)$．
\end{enumerate}
\end{theorem}
\begin{proof}
$A$ を含む有界閉区間 $I$ を固定する．

\textbf{(1)}\quad
$m^*(I \setminus A) \le m^*(I) = |I|$ だから
$m_*(A) = |I| - m^*(I \setminus A) \ge 0$．
$m_*(\emptyset) = |I| - m^*(I) = 0$．

\textbf{(2)}\quad
$I = A \sqcup (I \setminus A)$ であり，
可算劣加法性から
$|I| = m^*(I) \le m^*(A) + m^*(I \setminus A)$．
したがって
$m_*(A) = |I| - m^*(I \setminus A) \le m^*(A)$．

\textbf{(3)}\quad
$A \subset B$ ならば $I \setminus B \subset I \setminus A$，
よって $m^*(I \setminus B) \le m^*(I \setminus A)$ であり
$m_*(A) = |I| - m^*(I \setminus A) \le |I| - m^*(I \setminus B) = m_*(B)$．

\textbf{(4)}\quad 定義そのものである．

\textbf{(5)}\quad
$m^*$ の平行移動不変性から
$m^*((I+c) \setminus (A+c)) = m^*(I \setminus A)$，
$|I + c| = |I|$ であるから
$m_*(A+c) = |I+c| - m^*((I+c) \setminus (A+c))
= |I| - m^*(I \setminus A) = m_*(A)$．
\end{proof}

\begin{remark}
性質 (4) は Lebesgue内測度の定義を言い換えたものであるが，
測度論において本質的な役割を果たす：
集合 $A$ が可測であることの条件
「$m_*(A) = m^*(A)$」は，(4) により
「$m^*(A) + m^*(I \setminus A) = |I|$」と同値になる．
つまり $A$ とその補集合 $I \setminus A$ への分割で
外測度の加法性が成り立つことが可測性の本質である．
\end{remark}

\begin{theorem}[区間・基本集合の内測度]\label{thm:inner-elementary}
\begin{enumerate}
\item 半開区間 $Q$ に対して $m_*(Q) = |Q|$．
\item 基本集合 $E$ に対して $m_*(E) = m_J(E)$．
\end{enumerate}
\end{theorem}
\begin{proof}
$I \supset E$ を有界閉区間とする．
$I \setminus E$ は基本集合であり，
系\ref{cor:outer-measure-elementary}より
$m^*(I \setminus E) = m_J(I \setminus E) = |I| - m_J(E)$．
したがって
$m_*(E) = |I| - m^*(I \setminus E) = m_J(E)$．
半開区間は基本集合の特別な場合である．
\end{proof}

% ==================================================================
\section{具体例}\label{sec:lebesgue-inner-examples}
% ==================================================================

\begin{example}[{$\QQ \cap [0,1]$ の内測度}]\label{ex:rationals-inner}
$I = [0,1]$ とする．
$I \setminus (\QQ \cap [0,1])$
は $[0,1]$ 内の無理数全体であり，
$\QQ \cap [0,1]$ は可算だから
$m^*(\QQ \cap [0,1]) = 0$（例\ref{ex:rationals-outer}）．
可算劣加法性より
\[
1 = m^*([0,1])
\le m^*(\QQ \cap [0,1]) + m^*([0,1] \setminus \QQ)
= 0 + m^*([0,1] \setminus \QQ),
\]
かつ $m^*([0,1] \setminus \QQ) \le m^*([0,1]) = 1$ だから
$m^*([0,1] \setminus \QQ) = 1$．

したがって
\[
m_*(\QQ \cap [0,1])
= |[0,1]| - m^*([0,1] \setminus \QQ)
= 1 - 1 = 0.
\]

外測度 $m^*(\QQ \cap [0,1]) = 0$ と一致するから，
$\QQ \cap [0,1]$ はLebesgue可測であり，
そのLebesgue測度は $0$ である（第5章で正式に定義する）．
\end{example}

\begin{example}[Cantor集合の内測度]\label{ex:cantor-inner}
Cantor集合 $C$ に対して
$m^*(C) = 0$（例\ref{ex:cantor-outer}）だから
$0 \le m_*(C) \le m^*(C) = 0$
より $m_*(C) = 0$．
したがって $C$ もLebesgue可測で測度 $0$ であることが予想される
（第5章で確認する）．
\end{example}

\begin{example}[区間の内測度]\label{ex:interval-inner}
有界区間 $I = [a,b]$ に対して
$m_*(I) = m^*(I) = b - a$．
実際，$m_*(I) \le m^*(I) = b-a$ であり，
$J = [a,b]$ を取れば
$m_*(I) = |J| - m^*(J \setminus I) = (b-a) - 0 = b-a$．
\end{example}

\begin{example}[Jordan非可測だがLebesgue可測な集合]\label{ex:jordan-no-lebesgue-yes}
$\QQ \cap [0,1]$ は
Jordan内測度 $m_{J,*}(\QQ \cap [0,1]) = 0$ と
Jordan外測度 $m_J^*(\QQ \cap [0,1]) = 1$ が一致しないので
Jordan非可測であった．
しかし上で見たように
$m_*(\QQ \cap [0,1]) = m^*(\QQ \cap [0,1]) = 0$
だからLebesgue可測である．
Lebesgue理論はJordan理論より厳密に広いクラスの集合を扱える．
\end{example}

% ==================================================================
\section{4つの測度の大小関係}\label{sec:four-measures}
% ==================================================================

Jordan内測度・Lebesgue内測度・Lebesgue外測度・Jordan外測度の
間には，次の大小関係がある．

\begin{theorem}[4つの測度の大小関係]\label{thm:four-measures}
有界集合 $A \subset \RR^d$ に対して
\[
m_{J,*}(A) \le m_*(A) \le m^*(A) \le m_J^*(A).
\]
\end{theorem}
\begin{proof}
$m^*(A) \le m_J^*(A)$ は定理\ref{thm:jordan-lebesgue-outer}で示した．
$m_*(A) \le m^*(A)$ は定理\ref{thm:lebesgue-inner-properties}\,(2) で示した．
$m_{J,*}(A) \le m_*(A)$ を示す．

$I \supset A$ を有界閉区間，
$E \subset A$ を任意の基本集合とする．
$I \setminus A \subset I \setminus E$ であり
$I \setminus E$ は基本集合だから
\[
m^*(I \setminus A)
\le m^*(I \setminus E)
= m_J(I \setminus E)
= |I| - m_J(E).
\]
したがって
\[
m_*(A) = |I| - m^*(I \setminus A)
\ge |I| - (|I| - m_J(E))
= m_J(E).
\]
$E \subset A$ なる基本集合 $E$ について上限を取れば
$m_*(A) \ge m_{J,*}(A)$．
\end{proof}

\begin{corollary}\label{cor:jordan-implies-lebesgue}
$A$ がJordan可測（$m_{J,*}(A) = m_J^*(A)$）ならば，
$A$ はLebesgue可測（$m_*(A) = m^*(A)$）であり，
$m(A) = m_J(A)$．
\end{corollary}
\begin{proof}
$m_{J,*}(A) = m_J^*(A) =: c$ とすれば
$c \le m_*(A) \le m^*(A) \le c$
だから $m_*(A) = m^*(A) = c = m_J(A)$．
\end{proof}

\begin{remark}[解釈]
4つの測度の大小関係
\[
m_{J,*}(A) \le m_*(A) \le m^*(A) \le m_J^*(A)
\]
は，次のように解釈できる：
\begin{itemize}
\item 内測度は「内側からの評価」，外測度は「外側からの評価」であり，
常に内 $\le$ 外である．
\item Lebesgue版は可算被覆を用いるため，Jordan版より精密な評価を与える：
外測度は小さく（上からの評価が改善），
内測度は大きくなる（下からの評価が改善）．
\item Jordan可測ならば4つの値がすべて一致する．
\item 「Jordan可測でないがLebesgue可測」な集合
（例：$\QQ \cap [0,1]$）では，
中央の2つ（Lebesgue内外）は一致するが，
外側の2つ（Jordan内外）は一致しない．
\end{itemize}

\begin{center}
\begin{tabular}{l|cccc}
& $m_{J,*}$ & $m_*$ & $m^*$ & $m_J^*$ \\ \hline
$[0,1]$            & 1 & 1 & 1 & 1 \\
$\QQ \cap [0,1]$   & 0 & 0 & 0 & 1 \\
Cantor集合 $C$      & 0 & 0 & 0 & 0
\end{tabular}
\end{center}
\end{remark}
